% Owner: Member C (Eesha Irfan).
% Local font encoding restores the template's Times regular/bold faces under XeTeX.
\begingroup
\renewcommand{\encodingdefault}{T1}
\fontencoding{T1}\selectfont
\chapter{High-Level and Low-Level Design}
This chapter specifies the architectural framework, subsystem decomposition, domain models, and operational policies for Zeest. Designing an agentic medical copilot requires reconciling complex, non-deterministic foundation model reasoning with the stringent safety, traceability, and latency demands of clinical practice. The architecture transforms raw consultation audio, longitudinal electronic health records, and multimodal diagnostic inputs into structured case representations, evidence-backed differential diagnoses, and reviewable treatment plans. By decoupling specialized analytical agents from direct database writes, the system ensures that generative outputs remain provisional recommendations until explicitly validated by a licensed physician. The design bridges the functional specifications and data relationships established in Chapter~4 with concrete software engineering patterns, service boundaries, and state machines. Both high-level structural blueprints and low-level component contracts are formalized to provide a rigorous, transparent foundation for subsequent prototype realization. Every architectural component reflects strict adherence to clinical autonomy, ensuring automated assistance empowers rather than supplants human medical judgment across diverse consultation environments.

\textcolor{red}{PLANNED: The software architecture, agent topology, domain classes, and state machines described herein represent the proposed system design. Prototype implementation, hardware infrastructure, and empirical evaluation remain planned work.}

\section{System Overview}
Zeest is organized as a multi-tier, physician-governed clinical decision support system designed to assist practitioners during complex patient consultations. At its core, the software operates as an intelligent copilot that ingests streaming ambient audio, extracts clinically relevant concepts, retrieves longitudinal patient context, and generates transparent, cited diagnostic proposals. Rather than functioning as a monolithic artificial intelligence model, the system decomposes clinical problem solving into eight specialized agent modules coordinated through an explicit state graph. The presentation layer provides physicians with a responsive workspace for inspecting history, observing real-time transcription, reviewing multimodal diagnostic findings, and evaluating generated clinical drafts alongside supporting medical literature citations. Behind this interface, an asynchronous application gateway validates incoming requests, enforces authentication boundaries, and dispatches clinical workflows across dedicated machine learning inference runtimes and cloud-hosted foundation model application programming interfaces. This modular separation isolates analytical complexity while delivering a unified, intuitive clinical interface that preserves practitioner attention.

The overarching design is governed by the principle of human-in-the-loop clinical autonomy, ensuring that artificial intelligence serves strictly as an advisory mechanism. Longitudinal clinical context is maintained across multiple encounters through a hybrid memory architecture that pairs structured relational records with dense semantic vector embeddings. When a physician opens an encounter, the system constructs a consolidated clinical case representation by combining active consultation notes, laboratory findings, radiological image classifications, and relevant past medical history. This case context is evaluated by reasoning and critic agents to produce differential diagnosis proposals accompanied by passage-level clinical guideline citations. The physician retains absolute decision authority, choosing whether to approve, modify, or reject each draft revision. Every physician evaluation action is permanently preserved in an append-only audit trail, ensuring institutional accountability and clinical safety throughout the consultation lifecycle. By strictly barring unapproved machine hypotheses from entering the permanent health record, the architecture prevents speculative errors from propagating into future care episodes.

\section{Design Considerations}
Devising a dependable clinical decision support architecture necessitates a thorough evaluation of technical dependencies, regulatory boundaries, and software quality objectives. Clinical environments present unique operational challenges, including real-time consultation time constraints, catastrophic consequences of diagnostic hallucinations, and strict patient data confidentiality mandates. The architectural strategy must balance computational agility with defensive safety verifications, ensuring that automated recommendations remain transparent, auditable, and grounded in validated medical evidence. Designing for healthcare practitioners requires minimizing cognitive overhead while making algorithmic uncertainty and evidence provenance instantly accessible. These considerations directly influence component partitioning, asynchronous task distribution, and fail-safe recovery protocols throughout the system. The following subsections outline the foundational assumptions, environmental constraints, guiding design principles, and engineering methodologies that govern the system design of Zeest. These systematic considerations establish the operational boundaries within which the high-level system architecture, specialized agent subsystems, domain models, and decision governance policies are formally formulated, rigorously reviewed, and iteratively executed.

\subsection{Assumptions and Dependencies}
The design of Zeest relies on several key software, hardware, and operational assumptions regarding its operating environment and external service availability. Operationally, it is assumed that the primary end-user is a qualified medical practitioner possessing the clinical expertise required to critically evaluate diagnostic drafts, interpret radiological predictions, and verify pharmaceutical dosages. Technically, the system assumes stable network connectivity between the physician web client, the local application backend, and external model inference endpoints. The architecture depends upon modern web browser support for secure WebSockets and media capture interfaces to enable low-latency consultation audio streaming. Furthermore, the design assumes the availability of validated pre-trained neural networks for radiological screening and reliable medical knowledge bases for pharmacological contraindication checks, treating these components as modular service dependencies rather than training them from raw data. The persistence layer assumes relational transactional support and efficient vector similarity operations provided by an enterprise relational database engine equipped with high-dimensional indexing extensions.

\subsection{General Constraints}
The system architecture must operate within rigid computational, clinical safety, ethical, and regulatory constraints. Strict ethical boundaries prohibit the autonomous recording of clinical decisions or direct prescribing of medications without affirmative human confirmation. All generated diagnostic suggestions and treatment proposals must remain strictly segregated from immutable clinical health records until explicitly approved by the attending physician. Memory and storage constraints require efficient vector indexing techniques capable of scaling to extensive longitudinal patient histories without causing unacceptable query latency during live consultations. Data governance mandates require complete compliance with healthcare privacy regulations, dictating that all patient identifiers undergo synthetic pseudonymization prior to external model transit. Furthermore, hardware constraints necessitate that local compute burdens remain modest, offloading computationally intensive multimodal reasoning to scalable application programming interfaces while keeping local execution focused on lightweight orchestration and deterministic data validation. The system must also guarantee complete traceability, ensuring every generated recommendation remains linked to the exact input data, prompt configuration, and guideline passages that produced it.

\subsection{Goals and Guidelines}
The architectural design of Zeest is guided by core engineering principles emphasizing safety, transparency, modularity, and defensive failure isolation. Foremost among these is the principle of evidence groundedness, which dictates that no diagnostic hypothesis or therapeutic recommendation shall be presented without linked citations to verified medical literature or clinical guidelines. The user experience adheres to cognitive clarity guidelines, ensuring that provisional machine drafts are visually distinct from finalized clinical decisions to avoid user confusion. A strict modular separation of concerns is maintained across all specialist agents, enabling independent updates to individual diagnostic classifiers, embedding models, or reasoning prompts without destabilizing the broader orchestration pipeline. Finally, the system enforces defensive failure handling; if a specialized diagnostic agent or external guideline retrieval service encounters a failure, the orchestrator gracefully degrades performance, displaying transparent status indicators rather than producing unverified diagnostic speculations. These guidelines collectively ensure that the resulting system provides robust clinical utility while upholding the highest standards of software engineering craftsmanship and medical ethics.

\subsection{ Development Methods}
The development and specification of the Zeest architecture follow an iterative, specification-driven software engineering methodology adapted for complex multi-agent artificial intelligence systems. Architectural requirements and data models were formally derived from the functional requirements, quality attributes, and entity relationships established in Chapter~4. Component boundaries were established using standard object-oriented design patterns, domain-driven design principles, and unified modeling language structural and behavioral notations. Given the rapid evolution of foundation model capabilities and agentic orchestration frameworks, an agile prototyping strategy is adopted. Interface contracts, data serialization schemas, and error boundaries are rigorously formalized before implementing underlying model integrations. This structured methodology ensures complete traceability between high-level project objectives, clinical safety policies, and low-level code implementations, mitigating architectural drift throughout the engineering lifecycle. By establishing clear contracts and defensive abstractions early, the engineering team maintains the flexibility needed to evaluate candidate foundation models and specialized classifiers without rewriting core clinical business logic or compromising overarching system reliability.

\section{System Architecture}
The system architecture of Zeest is organized as a multi-tier, decoupled distributed system structured into three primary horizontal layers: the client and gateway layer, the orchestration and specialist agents layer, and the services and persistence layer. This tiered arrangement isolates user interface rendering, application routing, multi-agent reasoning, machine learning inference runtimes, and persistent data storage into distinct computational domains. By enforcing unidirectional dependencies across tiers, the architecture ensures that client interactions cannot bypass security controls, and non-deterministic agent workflows cannot directly mutate permanent health records. Layered decoupling also enables the team to scale computation-heavy inference services independently from the stateful application gateway and persistent relational storage backends. Figure~\ref{fig:esha:d07} illustrates the complete architectural topology, depicting inter-component communication channels, data flow boundaries, and external technology dependencies. The diagram shows how clinical inputs pass from the physician web client through the application gateway into the orchestration graph, which coordinates specialist analysis before returning evidence-backed proposals.

\begin{figure}[htbp]
\centering
\includegraphics[width=0.96\textwidth,height=0.82\textheight,keepaspectratio]{Figures/esha_D07_system_architecture.png}
\caption{Proposed high-level system architecture of Zeest (D07), illustrating the presentation tier, FastAPI application gateway, LangGraph multi-agent orchestration tier, specialist diagnostic agents, model inference runtimes, and PostgreSQL/pgvector persistence layer. Dashed lines denote internal component structures; solid lines represent cross-tier communication.}
\label{fig:esha:d07}
\end{figure}

The presentation tier is implemented as a modern web application constructed with Next.js, React, and TypeScript~\cite{zeestC:nextjs}. It provides specialized interface panels for consultation audio capture, structured note editing, multimodal input uploading, diagnostic review, and immutable audit inspection. The web client communicates with the application gateway via secure HTTPS RESTful endpoints for transactional case operations and full-duplex WebSockets for streaming consultation audio. The application gateway, built using FastAPI and Python~\cite{zeestC:fastapi}, serves as the secure entry point to the system. It encapsulates authentication guards implementing password hashing policies, validates incoming request payloads against strict schemas, and manages physician sessions. The gateway shields internal agent subsystems from direct client exposure, dispatching structured execution graphs to the orchestration tier while managing database transactions for review actions and decision recording. By enforcing strict request validation at the gateway boundary, invalid clinical inputs and malformed requests are rejected immediately, protecting internal reasoning workflows from inconsistent data states.

The core computational intelligence resides within the agent orchestration tier, which utilizes a state-graph engine based on LangGraph~\cite{zeestC:langgraph} and Google Agent Development Kit concepts. The orchestrator maintains an execution state context containing all intermediate findings, retrieved historical notes, and active draft revisions for the current encounter. It sequences the execution of eight specialist agents, coordinating parallel analyses and aggregating diagnostic outputs. The specialist agents interact with dedicated inference runtimes, including pre-trained PyTorch neural networks for chest radiograph classification~\cite{zeestC:pytorch}, speech recognition engines for audio transcription, and the Google Gemini foundation model application programming interface for clinical reasoning and safety critique~\cite{zeestC:gemini}. Finally, the persistence tier integrates a relational PostgreSQL database for atomic encounter records and audit trails with pgvector~\cite{zeestC:pgvector} for cosine similarity search over patient semantic histories and medical guideline embeddings. This architecture achieves robust isolation between fast relational operations and computationally demanding vector searches, providing high throughput and predictable query latency.

\subsection{Subsystem Architecture}
The multi-agent subsystem decomposes complex clinical diagnostic reasoning into eight modular, collaborating specialist agents. Each agent possesses a narrowly defined responsibility, a structured input-output contract, and dedicated tool access, preventing cross-domain contamination and facilitating targeted auditing. By decomposing clinical analysis into dedicated specialists, the system mirrors the multidisciplinary consultative process observed in institutional medical practice. Rather than forcing a single generic language model to parse images, calculate drug dosages, and recall obscure guidelines simultaneously, specialized modules handle specific data representations with dedicated tools. Figure~\ref{fig:esha:d08} illustrates the coordination topology, categorizing the agents into three functional pipeline stages: context and intake specialists, domain-specific diagnostic specialists, and synthesis, retrieval, and governance specialists. This structured progression ensures that patient history is retrieved and consolidated before domain-specific analyses begin, and that specialized findings are fully synthesized before clinical recommendations are drafted and critiqued. By standardizing message passing across these sequential boundaries, the orchestrator ensures that each specialist operates with maximum context while maintaining independent diagnostic reasoning channels.

\begin{figure}[htbp]
\centering
\includegraphics[width=0.96\textwidth,height=0.78\textheight,keepaspectratio]{Figures/esha_D08_agent_subsystems.png}
\caption{Proposed specialist agent subsystems and coordination topology (D08). The architecture organizes the eight agents into three sequential stages supervised by the multi-agent orchestration engine, enforcing clear input-output contracts and governance boundaries.}
\label{fig:esha:d08}
\end{figure}

Stage~1 comprises the Patient Memory Agent and the Clinical Case Agent. The Patient Memory Agent retrieves prior relational encounters and performs vector similarity search over historical semantic embeddings, reconstructing longitudinal patient context. The Clinical Case Agent aggregates these historical summaries with active consultation notes and newly submitted clinical inputs, synthesizing a unified case context. Stage~2 comprises the Imaging Agent, Laboratory Agent, and Medication Safety Agent. The Imaging Agent processes uploaded radiological scans through a pre-trained PyTorch classifier to detect pathological features. The Laboratory Agent analyzes structured metabolic panels, flagging abnormal values and identifying temporal biomarker trends. The Medication Safety Agent cross-references current prescriptions and proposed therapies against pharmacological knowledge bases, generating warnings for adverse drug interactions and contraindications. These specialized evaluations execute concurrently within the orchestration graph, maximizing computational throughput while maintaining independent diagnostic reasoning channels across distinct medical input modalities and computational runtimes. Each specialist publishes its structured findings directly into the shared working memory cache, making intermediate results immediately accessible to downstream reasoning and governance components.

Stage~3 encompasses the Evidence and RAG Agent, the Clinical Reasoning Agent, and the Safety and Critic Agent. The Evidence and RAG Agent queries guideline vector stores to retrieve verified clinical literature supporting active diagnostic hypotheses. The Clinical Reasoning Agent synthesizes the unified case context, specialist findings, and retrieved literature, invoking the Gemini foundation model to draft a structured differential diagnosis and treatment recommendation. Before presentation to the physician, the draft is audited by the Safety and Critic Agent, which inspects the proposal for logical consistency, unsupported clinical claims, and guideline divergence. This multi-layered synthesis ensures that generative inferences are thoroughly grounded in empirical patient findings and published medical evidence before reaching the user. Table~\ref{tab:esha:contracts} formalizes the operational contracts governing each specialist agent, defining their primary inputs, outputs, persistence mechanisms, and recipient components within the orchestration pipeline. These contracts establish clear data interfaces that guarantee modular maintainability and facilitate independent testing of individual agent subsystems.

\begin{table}[htbp]
\centering
\small
\caption{Proposed specialist agent input, output, persistence and recipient contracts (T14).}
\label{tab:esha:contracts}
\begin{tabular}{|p{0.18\textwidth}|p{0.22\textwidth}|p{0.25\textwidth}|p{0.25\textwidth}|}
\hline
Agent Name & Primary Inputs & Primary Outputs & Storage \& Recipient \\ \hline
Patient Memory Agent & \texttt{patient\_id}, encounter context & Relational history, top-$k$ semantic memories & PostgreSQL/pgvector; Orchestrator \\ \hline
Clinical Case Agent & Notes, multimodal data, retrieved history & Consolidated case representation & Working memory cache; Specialist agents \\ \hline
Imaging Agent & Radiological image tensor / file reference & Predicted pathology classes, confidence scores & Working cache; Reasoning Agent \\ \hline
Laboratory Agent & Laboratory panel values, reference ranges & Abnormal biomarker flags, trajectory trends & Working cache; Reasoning Agent \\ \hline
Medication Safety Agent & Active prescriptions, candidate drugs & Drug-drug interactions, contraindication alerts & Knowledge base; Reasoning Agent \\ \hline
Evidence \& RAG Agent & Diagnostic keywords, clinical hypotheses & Cited guideline passages, source identifiers & pgvector vector index; Reasoning Agent \\ \hline
Clinical Reasoning Agent & Consolidated case, specialist findings, evidence & Differential diagnosis draft, treatment plan & Gemini API; Safety \& Critic Agent \\ \hline
Safety \& Critic Agent & Generated draft, case context, cited evidence & Consistency evaluation, risk alerts, critique & Working cache; Presentation tier \\ \hline
\end{tabular}
\end{table}

The operational dynamics of clinical recommendation generation are depicted in Figure~\ref{fig:esha:d10}. The sequence begins when the physician initiates ambient audio capture during an encounter. Audio packets stream to the gateway via WebSockets, where an automatic speech recognition service converts the dialogue into structured consultation notes. The physician reviews, edits, and saves these notes, optionally uploading supplementary radiological images, laboratory values, or medication context. Upon the physician's explicit request for clinical analysis, the gateway triggers the orchestration workflow. The orchestrator queries the Patient Memory Agent, synthesizes case context, executes parallel diagnostic analyses, and retrieves supporting guideline citations. The Clinical Reasoning Agent drafts a differential diagnosis, which is evaluated by the Safety and Critic Agent. Once critique verification concludes, the draft proposal is saved with a pending review status, and the complete recommendation package is returned to the web interface for physician appraisal. This structured four-phase workflow guarantees that evidence is thoroughly cross-examined before entering the physician's cognitive workspace.

\begin{figure}[htbp]
\centering
\includegraphics[width=0.96\textwidth,height=0.78\textheight,keepaspectratio]{Figures/esha_D10_consultation_sequence.png}
\caption{Proposed clinical consultation and recommendation generation sequence (D10). The message sequence illustrates four operational phases: consultation capture, context retrieval, specialist analysis, and reasoning synthesis with safety verification.}
\label{fig:esha:d10}
\end{figure}

The decision review and memory persistence lifecycle is modeled in Figure~\ref{fig:esha:d12}. When reviewing a pending proposal revision, the physician may modify the draft text, creating a new revision while preserving the prior revision in the audit history. When the physician explicitly approves a draft revision, the gateway initiates an atomic database transaction. Within this transaction, an approval action is recorded in the audit log, the proposal status transitions to approved, and an immutable clinical decision record is persisted. Following transaction commitment, an asynchronous indexing worker generates dense text embeddings of the approved decision, storing the vector in pgvector. Conversely, if the physician rejects a proposal, the rejection action and status update are committed to the audit trail, but the content is strictly excluded from clinical decision and semantic memory tables. This atomic transactional boundary ensures that only verified clinical truths enter the longitudinal patient record, preventing unapproved machine hypotheses from influencing subsequent consultations.

\begin{figure}[htbp]
\centering
\includegraphics[width=0.96\textwidth,height=0.78\textheight,keepaspectratio]{Figures/esha_D12_memory_update_sequence.png}
\caption{Proposed physician review, audit logging, and memory update sequence (D12). The sequence demonstrates the transactional boundary governing approval, the audit trail of modifications and rejections, and the strict exclusion of unapproved drafts from semantic memory.}
\label{fig:esha:d12}
\end{figure}

\section{Architectural Strategies}
The architecture of Zeest incorporates deliberate high-level strategies designed to balance diagnostic accuracy, responsiveness, auditability, and clinical governance. These strategies resolve core tensions inherent in applying generative artificial intelligence to healthcare, specifically mitigating model hallucination, managing unstructured multimodal data, and enforcing strict human oversight. In complex clinical workflows, naive architectures that pipe unconstrained model outputs directly into patient charts present intolerable medical and legal hazards. The system resolves these challenges through two foundational architectural strategies: a hybrid memory persistence model that harmonizes transactional relational integrity with dense semantic search, and an asynchronous orchestration model that separates deterministic application logic from remote model runtimes. Architectural strategies also guide how network errors, partial computational failures, and token limits are managed without disrupting user interactions. Table~\ref{tab:esha:decisions} summarizes the key architectural dimensions, comparing supported design choices against evaluated alternatives and identifying unresolved technical parameters. By grounding high-level strategies in explicit engineering trade-offs, the design establishes transparent technical rationale that supports ongoing research exploration while preserving system maintainability.

\subsection{Hybrid Relational and Vector Memory Persistence}
Zeest adopts a dual-layer persistence strategy combining an ACID-compliant relational database with dense vector similarity indexing, implemented via PostgreSQL and pgvector. Clinical operations require absolute relational integrity for patient identity management, encounter tracking, draft versioning, and legal audit trails. Relational foreign keys ensure that every recorded action is inextricably bound to a specific physician, encounter, and proposal revision. However, relational schemas cannot effectively capture subtle semantic similarities across unstructured clinical narratives spanning years of patient care. To resolve this limitation, explicitly approved clinical decisions and intake notes are vectorized and indexed within pgvector. During subsequent consultations, semantic similarity queries retrieve clinically relevant historical episodes that keyword searches would miss, while relational joins enforce temporal ordering and patient isolation. This hybrid arrangement provides the mathematical power of semantic vector retrieval without compromising the transactional rigor, data integrity, and compliance guarantees demanded by modern healthcare information systems. By coupling pgvector similarity operators directly with PostgreSQL relational tables, the architecture achieves unified data administration without requiring independent synchronization mechanisms between disparate database engines.

\subsection{Asynchronous Agent Orchestration and Model Runtime Segregation}
To maintain interactive responsiveness during clinical encounters, the architecture employs an asynchronous orchestration strategy that physically segregates deterministic application logic from specialized model inference runtimes. Coordinating multiple artificial intelligence specialists sequentially would introduce unacceptable latency into the physician consultation workflow. Zeest addresses this by utilizing LangGraph state graphs to execute independent diagnostic analyses concurrently. The Imaging Agent, Laboratory Agent, and Evidence and RAG Agent operate in parallel threads once case context is established. Furthermore, external foundation model integrations are isolated behind strongly typed client adapters, preventing network timeouts or remote API errors from crashing the core application gateway. Fallback policies ensure that if an analytical service fails, the system returns partial, explicitly flagged findings rather than stalling the entire clinical workflow. This decoupled design enables independent horizontal scaling of inference microservices while keeping the application gateway lightweight, responsive, and resilient against third-party service interruptions or quota throttling. Furthermore, encapsulating foundation model interactions within asynchronous service boundaries prevents latency spikes during cloud model generation from blocking essential user interface operations or ambient consultation recording.

\begin{table}[htbp]
\centering
\small
\caption{Architectural design decisions, supported choices, trade-offs and unresolved options (T15).}
\label{tab:esha:decisions}
\begin{tabular}{|p{0.20\textwidth}|p{0.22\textwidth}|p{0.25\textwidth}|p{0.23\textwidth}|}
\hline
Design Dimension & Proposed Choice & Architectural Trade-off & Unresolved Status \\ \hline
Backend Framework & FastAPI (Python) & High-speed async I/O and native Python ML ecosystem versus separate microservices & Production server concurrency tuning TBD \\ \hline
Frontend Framework & Next.js / TypeScript & Rich React component ecosystem and streaming support versus heavier bundle size & View layout refactoring planned \\ \hline
Multi-Agent Orchestrator & LangGraph / Google ADK StateGraph & Cyclic state-graph flexibility versus increased execution graph complexity & Formal division of roles between frameworks TBD \\ \hline
Primary Reasoning Model & Google Gemini API & Advanced multimodal context window versus reliance on cloud API latency & Production model version and quotas TBD \\ \hline
Radiology Classifier & PyTorch Pre-trained Model & Validated local inference speed versus fixed diagnostic pathology label sets & Exact pre-trained model checkpoint TBD \\ \hline
Memory Storage & PostgreSQL with pgvector extension & Unified relational ACID guarantees and vector search versus dedicated vector database & Vector dimension and HNSW index tuning TBD \\ \hline
\end{tabular}
\end{table}

\section{Domain Model/Class Diagram}
The structural design of the Zeest domain layer maps the entity relationships and data dictionary definitions established in Chapter~4 into cohesive object-oriented software abstractions. The domain model separates core clinical entities, recommendation and governance structures, and domain service controllers into strongly encapsulated packages. This encapsulation ensures that clinical entities maintain domain invariants, lifecycle state machines are strictly enforced, and cross-cutting persistence operations are managed through dedicated service interfaces. By organizing domain logic into distinct object hierarchies, the architecture promotes high cohesion and low coupling across analytical modules. Clear separation of responsibilities also ensures that changes to underlying machine learning model interfaces do not require alterations to core patient accounting or encounter management classes. Figure~\ref{fig:esha:d09} presents the unified class diagram illustrating attributes, primary operations, and structural associations across all domain entities, value objects, and service coordinators. By translating the data dictionary and requirements of Chapter~4 into strongly typed class definitions, the domain architecture establishes a cohesive object model that directly supports subsequent software implementation and unit testing.

\begin{figure}[htbp]
\centering
\includegraphics[width=0.96\textwidth,height=0.82\textheight,keepaspectratio]{Figures/esha_D09_domain_classes.png}
\caption{Proposed domain model and service class diagram (D09). The diagram illustrates core clinical entities, recommendation and governance classes, enumerated lifecycle states, and domain service interfaces implementing the Chapter~4 data dictionary.}
\label{fig:esha:d09}
\end{figure}

The core entity package models foundational healthcare concepts: \texttt{Physician}, \texttt{Patient}, \texttt{Encounter}, and \texttt{ClinicalInput}. The \texttt{Physician} entity encapsulates credential verification and session management. A \texttt{Patient} maintains a collection of historical encounters, providing methods to synthesize longitudinal clinical summaries. Each \texttt{Encounter} aggregates consultation notes and multimodal findings represented as \texttt{ClinicalInput} records, which support versioned text corrections and binary file references. The recommendation and governance package models the diagnostic lifecycle. The \texttt{Proposal} entity tracks generated diagnostic text, critic feedback, integer revision numbers, and pending review states. Associated \texttt{EvidenceItem} entities capture cited medical passages and bibliographic references, linked to specific proposal revisions through \texttt{ProposalEvidence} associations. These entities ensure that every machine-generated recommendation remains linked to its supporting data, enabling practitioners to trace claims back to original clinical inputs and published medical literature without ambiguity or loss of provenance. In particular, the Proposal entity encapsulates the collaborative interaction between clinical reasoning and safety critique, ensuring that diagnostic text and critic feedback remain versioned together across successive review iterations.

Accountability and decision boundaries are enforced through \texttt{PhysicianAction}, \texttt{ClinicalDecision}, and \texttt{SemanticMemory}. Every review interaction instantiates a \texttt{PhysicianAction} recording the action kind, the responsible physician identifier, a precise timestamp, and a snapshot of reviewed text. The \texttt{ClinicalDecision} entity represents the finalized, physician-accepted conclusion, requiring a direct reference to the corresponding approval action and revision number. The \texttt{SemanticMemory} entity encapsulates dense vector embeddings, maintaining conditional foreign keys to originating inputs or accepted decisions. A dedicated domain service layer orchestrates business logic across entities: the \texttt{OrchestrationService} manages agent execution workflows, the \texttt{SafetyCriticService} validates diagnostic groundedness, the \texttt{DecisionService} executes transactional approvals, and the \texttt{MemoryService} performs cosine similarity retrieval over pgvector stores. This service architecture isolates complex business logic from data storage mechanisms, ensuring that changes to indexing algorithms or storage backends do not disrupt clinical workflows or compromise transaction boundaries. Furthermore, by encapsulating transactional decision recording within the DecisionService, the system guarantees that approval actions, status transitions, and clinical decision insertions execute within an atomic database boundary.

\section{Policies and Tactics}
Policies and tactics govern localized implementation decisions, interface rules, and defensive software patterns that do not alter top-level structural decomposition but directly impact system reliability and clinical safety. In Zeest, these mechanisms enforce rigid human oversight, guarantee legal auditability, and eliminate unchecked foundation model hallucinations. In safety-critical healthcare software, robust architectural decomposition alone is insufficient; specific computational policies must govern state transitions, decision recording, and output verification. These defensive policies establish automated checkpoints throughout the consultation lifecycle, ensuring that algorithmic suggestions are systematically verified before reaching practitioner attention. By embedding programmatic safety guards into core service protocols, the architecture minimizes the likelihood of human error while safeguarding medical record integrity. The subsections below delineate the governing approval policy and the evidence citation verification tactic, detailing the concrete mechanisms through which clinical safety and institutional accountability are preserved. These policies translate ethical and clinical guidelines into programmatic constraints, ensuring that software behaviors strictly conform to the human-in-the-loop mandate across all operational scenarios.

\subsection{Physician Review Boundary and Audit Immutability Policy}
The central operational policy of Zeest mandates an inviolable boundary between provisional artificial intelligence proposals and authoritative clinical records. Under no circumstances may an automated agent instantiate a \texttt{ClinicalDecision} record. As modeled in the state machine of Figure~\ref{fig:esha:d11}, every generated draft enters a mandatory \texttt{PendingReview} state. When a physician edits a draft, the system transitions to \texttt{ModifiedDraftPending}, creating an incremented revision while preserving the initial draft in an immutable audit entry. Only an explicit approval action by a validated physician can transition a proposal to \texttt{ExplicitlyApproved}, atomically writing the accepted content to the decision ledger. If a proposal is rejected, it transitions to \texttt{ExplicitlyRejected}; the rejection is permanently logged in the audit trail for institutional review, but the proposal is strictly barred from clinical records and semantic memory stores. This strict state transition policy guarantees that the system maintains an uncompromised audit trail of every interaction, safeguarding practitioner accountability while preventing speculative machine hypotheses from contaminating longitudinal patient histories.

\begin{figure}[htbp]
\centering
\includegraphics[width=0.88\textwidth,height=0.76\textheight,keepaspectratio]{Figures/esha_D11_approval_states.png}
\caption{Proposed clinical proposal lifecycle and decision state transitions (D11). The state machine models the progression from drafting to pending review, revision modifications, explicit approval, semantic memory indexing, and audit logging of rejections.}
\label{fig:esha:d11}
\end{figure}

\subsection{Grounded Clinical Evidence Citation and Critic Verification Tactic}
To mitigate generative hallucinations and ensure clinical trustworthiness, Zeest implements a strict evidence citation and safety verification tactic. When the Clinical Reasoning Agent produces a diagnostic recommendation, it is programmatically required to anchor each clinical claim to supporting passage identifiers supplied by the Evidence and RAG Agent. Prior to presenting the draft on the physician dashboard, the Safety and Critic Agent evaluates the proposal against the source evidence. The critic verifies that cited passages genuinely substantiate the claims, checks for contraindications flagged by the Medication Safety Agent, and inspects laboratory value consistency. If unsupported assertions, conflicting diagnoses, or ungrounded claims are identified, the critic appends structured warning flags to the proposal payload. These warnings are rendered prominently on the physician review interface, providing practitioners with critical context before they decide to approve, modify, or reject the recommendation. This two-tier verification tactic ensures that physicians receive transparent, fact-checked diagnostic guidance rather than opaque foundation model completions.

\endgroup
